\subsection{Two limitations of a direct eleventh-Kummer comparison}

Put $D=\ker\lambda\cap\bigcap_q\ker\rho_q^I$ in the full
profinite Grothendieck--Teichm\"uller group.  By \cite[Theorem~8.1 and Remark~8.6]{Enriquez2007}, the exact
Kummer relation, with zero Kummer coordinate on $D$, is
\begin{equation}\label{eq:eleventh-local-system-K}
d_n(f_g)=f_g\qquad(g\in D).
\end{equation}
Here $H_n\leq\widehat F(x,y)$ is the kernel of $x\mapsto1,y\mapsto0$
in $\mathbb Z/n$, freely generated by
\[
 X=x^n,\qquad Y_j=x^{-j}yx^j\quad(0\leq j<n),
\]
and $d_n:H_n\to\widehat F(x,y)$ sends $X\mapsto x$, $Y_0\mapsto y$
and $Y_j\mapsto1$ for $j\ne0$.  We apply the full nonabelian relation \eqref{eq:eleventh-local-system-K}.

\begin{lemma}\label{lem:kn-different-local-systems}
For $n=11$ and a primitive cube root $\omega$ in
$\mathbb F_{17^2}$, consider characters over $\overline{\mathbb F}_{17}$.
The character of $H_{11}$ obtained by pulling
back $x\mapsto1,y\mapsto\omega$ along $d_{11}$ is different from
the restriction of every character
$x\mapsto\xi,y\mapsto\omega$ with $\xi^{11}=1$.
Nor is either kind a deck translate of the other.
\end{lemma}
\begin{proof}
The pullback along $d_{11}$ has values
\[
 X\mapsto1,\qquad Y_0\mapsto\omega,\qquad
 Y_j\mapsto1\ (j\ne0).
\]
Each character of the second kind has values
\[
 X\mapsto\xi^{11}=1,\qquad Y_j\mapsto\omega
 \quad\hbox{for every }j.
\]
The deck generator cyclically permutes the conjugacy classes of
the $Y_j$.  It therefore permutes the eleven characters having a
single nontrivial $Y_j$-value and fixes the character whose eleven
$Y_j$-values are all $\omega$.  These two orbits are disjoint.
\end{proof}

Consequently the direct operation ``apply \eqref{eq:eleventh-local-system-K}, then pass to the
pulled-back rank-one local system'' does not identify the product
\begin{equation}\label{eq:eleventh-local-system-P}
P_g=\prod_{\substack{\xi^{11}=1\\\xi\ne1}}B_g(\xi,\omega)
\end{equation}
with the determinant on that pulled-back local system.  This
observation is a limitation of this proposed argument, not a claim
that the full relation \eqref{eq:eleventh-local-system-K} cannot constrain \eqref{eq:eleventh-local-system-P} by another route.

Here is a separate, explicit limitation of a pairing-only argument.
Let $C$ be the smooth compact curve with affine equation
$u^{11}+v^3=1$, and let
$A=C_{11}\times C_3$ be its deck group over the $t$-line,
$t=u^{11}$.  We use the geometric finite homology description:
over an algebraic closure of $\mathbb F_{17}$, its $20$-dimensional
homology has exactly the character lines
\begin{equation}\label{eq:eleventh-local-system-H}
(\xi^a,\omega^b),\qquad
 a\in(\mathbb Z/11)^\times,\quad b\in\{1,2\}.
\end{equation}
Every line is one-dimensional.  The alternating intersection form
pairs a line only with its inverse-character line.
For completeness, the character description follows from the
two-loop cellular complex of the three-punctured sphere.  On a
nontrivial character the boundary from the two one-cells to the
one zero-cell has rank one, so the open cover has a single
homology line.  Compactification kills precisely the peripheral
lines whose character is trivial on one of the three inertia
groups.  The inertia values here are $\xi^a$, $\omega^b$ and
$(\xi^a\omega^b)^{-1}$.  Since eleven and three are coprime,
all three are nontrivial exactly for the pairs in \eqref{eq:eleventh-local-system-H}.
The trivial character contributes no compact homology.

\begin{proposition}\label{prop:kn-linear-sign-model}
There is an $\mathbb F_{17}$-linear automorphism of the space \eqref{eq:eleventh-local-system-H}
that commutes with the whole deck group $A$, preserves its
alternating intersection form, and has determinant $-1$ on the
ten-dimensional $\omega$-eigenspace over $\mathbb F_{17^2}$.
Its quadratic norm-one determinant sign is therefore $-1$.
This is a linear model, not a GT element or a solution of the
pentagon or of \eqref{eq:eleventh-local-system-K}.
\end{proposition}
\begin{proof}
The square of Frobenius acts on the nonzero exponents modulo
eleven by multiplication by $17^2\equiv3$.  Its two orbits are
\[
 A_0=\{1,3,9,5,4\},\qquad
 A_1=\{2,6,7,10,8\}=-A_0.
\]
Ordinary Frobenius acts on the pairs $(a,b)$ in \eqref{eq:eleventh-local-system-H} as
\[
 (a,b)\longmapsto(6a,-b).
\]
It thus preserves the set
\[
 \mathcal O=(A_0\times\{1\})\cup(A_1\times\{2\}).
\]
Indeed $6A_0=A_1$ and $6A_1=A_0$.
The set $\mathcal O$ is also stable under inversion
$(a,b)\mapsto(-a,-b)$, as is its complement.

Define the automorphism to be $-1$ on the ten character lines in
$\mathcal O$ and $+1$ on the ten remaining lines.  Frobenius
stability makes it descend to $\mathbb F_{17}$, and its scalar
action on each character line makes it commute with $A$.
The intersection form pairs only mutually inverse characters:
if $v,w$ have characters $\chi,\psi$, invariance gives
$\langle v,w\rangle=\chi(a)\psi(a)\langle v,w\rangle$ for every
deck element $a$, which forces $\psi=\chi^{-1}$ whenever the
pairing is nonzero.  Our automorphism has equal signs on these
paired lines, so it preserves that form.

There are exactly five negative lines in the $\omega$-eigenspace,
namely those indexed by $A_0$.  Its determinant there is
$(-1)^5=-1$.  This value belongs to the norm-one group of
$\mathbb F_{17^2}^\times$, and raising it to
$(17+1)/2=9$ leaves $-1$, as asserted.
\end{proof}

The exact remaining question is whether the full nonabelian
relations, including \eqref{eq:eleventh-local-system-K} and the pentagon, exclude this sign in
the GT action.  Neither Lemma~\ref{lem:kn-different-local-systems}
nor Proposition~\ref{prop:kn-linear-sign-model} answers it.
In particular, no assertion about nonvanishing of
$\Theta_{17,11}|_{\widehat{\mathrm{GT}}}$ follows from these two statements.
